% GENERATED FILE. Edit canonical lessons, then run npm run generate:books.
% canonical-source-hash: fnv1a64:21e9a2441d6716e4
% canonical-lessons: SA-C22-hear, SA-C22-lead, SA-C22-play, SA-C22-burn, SA-S221-letter-dda

\chapter{Verbs of Mind and Motion}
\label{ch:sa-verbs-of-mind}
\hlchaptermodality{22}

\begin{chapteropening}
\textbf{By the end of this chapter:} \emph{I can say that someone hears, leads, plays and asks in Sanskrit, and name the root behind each.}

Complete SA-C22-burn and say all four verbs in order.
\end{chapteropening}

\section[śṛṇoti]{\sk{शृणोति} (śṛṇoti) — he, she or it hears}

\label{lesson:SA-C22-hear}



\begin{quote}
A short run of everyday verbs, one at a time.
\end{quote}

\subsection*{You'll want to know: \sk{शृणोति}}
\textbf{\sk{शृणोति}} (\emph{śṛṇoti}) — \textquotedblleft{}he, she or it hears\textquotedblright{}.

From \textbf{√\sk{श्रु}} (\emph{śru}), 'to hear' — the same root as Latin \emph{clueo} and Greek \emph{kleos}, 'fame', which is what gets heard about you. It gives \textbf{\sk{श्रुति}} (\emph{śruti}), 'that which is heard': the Vedas, as against \emph{smṛti}, 'that which is remembered'.

\begin{grammarlens}[title={What you've built}]
One more everyday action. Three follow, and each reuses the ones before.
\end{grammarlens}

\subsection*{Your turn}
\begin{itemize}
\raggedright
  \item \emph{Say it:} \emph{śṛṇoti}
  \item \emph{Say it:} it once more, slowly
  \item \emph{Recall:} say \emph{naraḥ}, then read \textbf{\sk{नारी}}
\end{itemize}

\begin{tcolorbox}[breakable,colback=teal!4,colframe=teal!35!black,title={Before you move on}]
What does \sk{शृणोति} mean? (\textquotedblleft{}he, she or it hears\textquotedblright{}.) Say it once more.
\end{tcolorbox}
\section[nayati]{\sk{नयति} (nayati) — he, she or it leads}

\label{lesson:SA-C22-lead}



\begin{quote}
Before the new one: what did \emph{śṛṇoti} mean?
\end{quote}

\subsection*{You'll want to know: \sk{नयति}}
\textbf{\sk{नयति}} (\emph{nayati}) — \textquotedblleft{}he, she or it leads\textquotedblright{}.

From \textbf{√\sk{नी}} (\emph{nī}), 'to lead'. It is inside \emph{netṛ}, 'leader', and \emph{naya}, 'polity' — the art of leading people, which Sanskrit treats as a subject with a name.

\begin{grammarlens}[title={What you've built}]
One more everyday action. Three follow, and each reuses the ones before.
\end{grammarlens}

\subsection*{Your turn}
\begin{itemize}
\raggedright
  \item \emph{Say it:} \emph{nayati}
  \item \emph{Say it:} it once more, slowly
  \item \emph{Say it:} it after \emph{śṛṇoti}, so the two sit together
  \item \emph{Recall:} read \textbf{\sk{बालः}}
\end{itemize}

\begin{tcolorbox}[breakable,colback=teal!4,colframe=teal!35!black,title={Before you move on}]
What does \sk{नयति} mean? (\textquotedblleft{}he, she or it leads\textquotedblright{}.) Say it once more.
\end{tcolorbox}
\section[krīḍati]{\sk{क्रीडति} (krīḍati) — he, she or it plays}

\label{lesson:SA-C22-play}



\begin{quote}
Before the new one: what did \emph{nayati} mean?
\end{quote}

\subsection*{You'll want to know: \sk{क्रीडति}}
\textbf{\sk{क्रीडति}} (\emph{krīḍati}) — \textquotedblleft{}he, she or it plays\textquotedblright{}.

From \textbf{√\sk{क्रीड्}} (\emph{krīḍ}). The noun \emph{krīḍā} covers a child's game, a sport, and divine play — \emph{līlā}'s plainer cousin — so the word does not separate the trivial from the cosmic.

\begin{grammarlens}[title={What you've built}]
One more everyday action. Three follow, and each reuses the ones before.
\end{grammarlens}

\subsection*{Your turn}
\begin{itemize}
\raggedright
  \item \emph{Say it:} \emph{krīḍati}
  \item \emph{Say it:} it once more, slowly
  \item \emph{Say it:} it after \emph{nayati}, so the two sit together
  \item \emph{Recall:} say \emph{vṛddhaḥ}
\end{itemize}

\begin{tcolorbox}[breakable,colback=teal!4,colframe=teal!35!black,title={Before you move on}]
What does \sk{क्रीडति} mean? (\textquotedblleft{}he, she or it plays\textquotedblright{}.) Say it once more.
\end{tcolorbox}
\section[dahati]{\sk{दहति} (dahati) — he, she or it burns}

\label{lesson:SA-C22-burn}



\begin{quote}
Before the new one: what did \emph{krīḍati} mean?
\end{quote}

\subsection*{You'll want to know: \sk{दहति}}
\textbf{\sk{दहति}} (\emph{dahati}) — \textquotedblleft{}he, she or it burns\textquotedblright{}.

From \textbf{√\sk{दह्}} (\emph{dah}), and its distant English cousin is \textbf{day} — by way of a Germanic word for the burning warmth of the sky, not the daylight itself.

\begin{grammarlens}[title={What you've built}]
That closes the run: four more actions, each met on its own.
\end{grammarlens}

\subsection*{Your turn}
\begin{itemize}
\raggedright
  \item \emph{Say it:} \emph{dahati}
  \item \emph{Say it:} it once more, slowly
  \item \emph{Say it:} it after \emph{krīḍati}, so the two sit together
  \item \emph{Recall:} read \textbf{\sk{जनः}}
\end{itemize}

\begin{tcolorbox}[breakable,colback=teal!4,colframe=teal!35!black,title={Before you move on}]
What does \sk{दहति} mean? (\textquotedblleft{}he, she or it burns\textquotedblright{}.) Say it once more.
\end{tcolorbox}
\section[ḍa]{\sk{ड} — one character, met inside words you already say}

\label{lesson:SA-S221-letter-dda}



\begin{quote}
Before the new one, the ones you have already met: \sk{ज} · \sk{श} · \sk{छ}. Say what each of them does.

One character this time. One only — and you have been saying it for pages without knowing which mark on the page it was.
\end{quote}

\subsection*{Script you'll notice: \sk{ड}}
\textbf{\sk{ड}} — \emph{ḍa}.

It is a \textbf{consonant}, and in this script a consonant is never bare: it comes with an \emph{a} already in it. So it is not \emph{ḍ}, it is \textbf{ḍa}.

What it is made of:

\begin{itemize}
\raggedright
  \item a descending stem joined to an upper-left loop and broad open lower bowl
  \item the top shirorekhā
\end{itemize}

You already say these, and every one of them has \sk{ड} somewhere inside it:

\begin{itemize}
\raggedright
  \item \textbf{\sk{क्रीडति}} \emph{krīḍati} — he, she or it plays
\end{itemize}

\subsection*{Writing: \sk{ड}}
\begin{itemize}
\raggedright
  \item \textbf{1.} start at the stem's headline junction, descend, turn left across the upper shoulder, sweep counterclockwise around the upper-left loop, then continue clockwise around the broad lower bowl to its open up-left tip without lifting
  \item \textbf{2.} lift and draw the top shirorekhā left-to-right
\end{itemize}

\textbf{Pen lifts: 1.} The pen comes up 1 times and no more.

\begin{quote}
verified two-stroke teaching form fitted to the bundled printed outline
\end{quote}

\begin{quote}
This is one attested teaching order and not a national standard — handwriting here is taught with school-to-school variation. Source: Opiaterein, ‘Deva-\sk{ड}-order.gif’, strokes 1–2, Wikimedia Commons, 10 May 2009.
\end{quote}

\subsection*{Your turn}
\begin{itemize}
\raggedright
  \item \emph{Look:} at this, and find \sk{ड} in it
\end{itemize}

\begin{quote}
\sk{क्रीडति}
\end{quote}

\begin{itemize}
\raggedright
  \item \emph{Trace it:} \sk{ड} three times, saying \emph{ḍa} as you finish each one
  \item \emph{Look:} back at any page of this chapter and find \sk{ड} once more
\end{itemize}

\begin{tcolorbox}[breakable,colback=teal!4,colframe=teal!35!black,title={Before you move on}]
Which character is this — \sk{ड}? What sound does it carry? (\emph{\textbf{ḍa}}.) Name one word you already say that contains it.
\end{tcolorbox}
